% Moment matching for the Laplace-kernel covariance.
% Standalone method note. Builds with: pdflatex moment_estimator.tex (twice, for refs).
\documentclass[11pt,a4paper]{article}

\usepackage{lmodern}        % scalable fonts; microtype's expansion needs them
\usepackage[T1]{fontenc}
\usepackage[utf8]{inputenc}
\usepackage[margin=2.6cm]{geometry}
\usepackage{amsmath,amssymb}
\usepackage{booktabs}
\usepackage{array}
\usepackage{listings}
\usepackage{xcolor}
\usepackage{microtype}
\usepackage[hidelinks]{hyperref}

\newcommand{\code}[1]{\texttt{#1}}
\newcommand{\yes}{\checkmark}

\definecolor{codekw}{rgb}{0.30,0.10,0.50}
\definecolor{codecm}{rgb}{0.35,0.45,0.35}
\definecolor{codest}{rgb}{0.60,0.20,0.10}
\lstdefinestyle{py}{
  language=Python,
  basicstyle=\ttfamily\footnotesize,
  keywordstyle=\color{codekw}\bfseries,
  commentstyle=\color{codecm}\itshape,
  stringstyle=\color{codest},
  showstringspaces=false,
  columns=fullflexible,
  breaklines=true,
  frame=single,
  rulecolor=\color{black!25},
  framesep=5pt,
  xleftmargin=6pt,
  numbers=none,
}
\lstset{style=py}

\title{Moment matching for the Laplace-kernel covariance}
\author{Tomasz Ku\'smierczyk}
\date{27 October 2026}

\begin{document}
\maketitle

\begin{abstract}
\noindent
A closed-form way to fit the covariance family of \emph{``Stop Marginalizing My Dreams: Model
Inversion via Laplace Kernel for Continual Learning''} (arXiv:2605.11804, Krukowski et al.),
which that paper fits by gradient descent. Three bands of the empirical second moment identify
the parameters outright, in one pass, with no optimiser. On flattened feature maps at the
paper's own scale it improves on their gradient fit by $0.47$ to $1.13$ nats per dimension
while fitting $17.5$--$20\times$ faster. It does not replace their whole method: it requires
the coordinate ordering to be supplied, which their gradient fit is meant to discover. What
follows gives the derivation, the exact relationship to their procedure, measured limits, and
a proposal for combining the two. The note is self-contained.
\end{abstract}

\section{The family}

For a covariance over $N$ variables the paper proposes
\begin{equation}
  \Sigma = \operatorname{diag}(d) + \operatorname{diag}(w)\, K(a)\, \operatorname{diag}(w),
  \qquad
  K(a)_{ij} = e^{-|a_i - a_j|}.
  \label{eq:family}
\end{equation}
Each variable $i$ receives a coordinate $a_i$ on a line, an amplitude $w_i$ and a diagonal
(nugget) term $d_i$: three parameters per variable, against the $O(N^2)$ of a dense covariance.
Correlation decays exponentially in coordinate distance, so variables close on the line are
strongly dependent and distant ones nearly independent.

The coordinates are \emph{learned}, which is the paper's distinctive move. Sorting by $a$
induces an ordering over the variables, so the model is intended to discover for itself which
variables belong near each other.

\paragraph{Why it is cheap.} Sort the coordinates and the model is a Gauss--Markov chain: an
AR(1) latent process observed with independent noise, equivalently an Ornstein--Uhlenbeck or
Mat\'ern-$\tfrac12$ process in state-space form. Its precision is tridiagonal, so quadratic
forms and log-determinants cost $O(N \log N)$ by prefix scan, or $O(N)$ by a Kalman recursion,
instead of $O(N^3)$. That chain structure is also what makes the fit below possible.

\section{The estimator}
\label{sec:estimator}

Write $\varphi_p = e^{-(a_p - a_{p-1})} \in (0,1)$ for the decay between neighbours along the
sorted order. For three consecutive coordinates, \eqref{eq:family} gives
\begin{align}
  \Sigma_{p,\,p+1}   &= w_p\, w_{p+1}\, \varphi_{p+1}, \label{eq:b1}\\
  \Sigma_{p+1,\,p+2} &= w_{p+1}\, w_{p+2}\, \varphi_{p+2}, \label{eq:b2}\\
  \Sigma_{p,\,p+2}   &= w_p\, w_{p+2}\, \varphi_{p+1}\varphi_{p+2}. \label{eq:b3}
\end{align}
Multiplying \eqref{eq:b1} by \eqref{eq:b2} and dividing by \eqref{eq:b3} cancels both outer
amplitudes and both decays, leaving a single amplitude:
\begin{equation}
  \boxed{\;
  w_{p+1}^{2}
  = \frac{\Sigma_{p,\,p+1}\; \Sigma_{p+1,\,p+2}}{\Sigma_{p,\,p+2}}
  \;}
  \label{eq:key}
\end{equation}
Substituting empirical moments for the model entries identifies $w^2$ directly. Everything else
follows in sequence, each step closed-form. With
\[
  s_0 = \mathbb{E}\!\left[x_p^2\right], \qquad
  s_1 = \mathbb{E}\!\left[x_p x_{p+1}\right], \qquad
  s_2 = \mathbb{E}\!\left[x_p x_{p+2}\right]
\]
the diagonal and first two off-diagonal bands, taken \emph{along the given order}:

\begin{center}
\begin{tabular}{@{}lll@{}}
\toprule
Step & Formula & From \\
\midrule
amplitudes  & $w^2_{p+1} = s_{1,p}\, s_{1,p+1} / s_{2,p}$        & bands 1 and 2 \\
decays      & $\varphi_p = s_{1,p} / (w_p w_{p+1})$              & band 1 \\
coordinates & $a = \operatorname{cumsum}(-\log \varphi)$          & the decays \\
nugget      & $d = s_0 - w^2$                                     & the diagonal \\
\bottomrule
\end{tabular}
\end{center}

\paragraph{Cost.} One pass over the data for the three bands, then $O(N)$ arithmetic, and
$O(N)$ memory --- the $N \times N$ matrix is never formed. There is nothing to tune: no
learning rate, no convergence criterion, no divergence.

\paragraph{Two details the three bands force.} The endpoint amplitudes $w^2_1$ and $w^2_N$ are
not identified by three bands, so they are copied inward from their neighbours. And $w^2$ is
clamped to $[0,\, 0.999\, s_0]$ so the implied nugget stays positive.

\paragraph{Correctness.} The fit is exact to floating point when the target genuinely lies in
the family, except at those two endpoints. Checked against a dense Cholesky factorisation at
$N = 16\,384$: log-determinant agreeing to $3.1\mathrm{e}{-10}$, quadratic form to
$3.1\mathrm{e}{-7}$, all eigenvalues positive.

\subsection{The whole estimator}

Listing~\ref{lst:fit} is the entire method, with the clamping kept, since the degenerate cases
are where a naive transcription breaks. Executed against the reference implementation on data
drawn from a known chain it reproduces it exactly, to $0.0$ difference on both the coordinates
and the amplitudes --- so this is runnable code, not pseudo-code.

\begin{lstlisting}[style=py,caption={The closed-form fit in full.},label={lst:fit},float=htbp]
def fit_moments(X, order=None, eps=1e-6, clip=0.999):
    """Closed-form LCM fit. X is (n_samples, dim), centred. Returns (a, w, d)
    in the ORIGINAL indexing, so the caller's variable order is preserved."""
    n, dim = X.shape
    perm = torch.arange(dim) if order is None else torch.as_tensor(order)
    P = X[:, perm]                                   # into chain order

    # the three bands -- the only statistics the fit ever sees
    s0 = (P * P).mean(0)                             # diagonal
    s1 = (P[:, :-1] * P[:, 1:]).mean(0)              # band 1
    s2 = (P[:, :-2] * P[:, 2:]).mean(0)              # band 2

    tol = eps * s0.mean().clamp(min=eps)

    # amplitudes: w^2_{p+1} = s1_p * s1_{p+1} / s2_p
    w2 = torch.zeros(dim, dtype=X.dtype)
    w2[1:-1] = torch.where(s2.abs() > tol, s1[:-1] * s1[1:] / s2, torch.zeros_like(s2))
    w2[0], w2[-1] = w2[1], w2[-2]                    # endpoints unidentified; copy inward
    w2 = w2.clamp(min=0.0).minimum(clip * s0)        # keep the nugget positive
    w = w2.sqrt()                                    # <-- sign is discarded here (Sec. 5)

    # decays from band 1, then coordinates as a cumulative sum of -log phi
    denom = (w[:-1] * w[1:]).clamp(min=tol)
    phi = (s1 / denom).clamp(min=tol, max=1.0 - 1e-12)
    a_chain = torch.cat([torch.zeros(1, dtype=X.dtype), (-phi.log()).cumsum(0)])

    d = (s0 - w2).clamp(min=tol)                     # nugget

    a, wv, dv = (torch.empty(dim, dtype=X.dtype) for _ in range(3))
    a[perm], wv[perm], dv[perm] = a_chain, w, d
    return a, wv, dv
\end{lstlisting}

Three lines carry the whole idea: the \code{w2[1:-1]} ratio, the \code{phi} division, and the
\code{cumsum}. Everything else is bookkeeping and guards. Note that $a$ is never initialised
and never updated --- it is \emph{computed}, which is the difference from the gradient fit
discussed in Section~\ref{sec:relation}. In library form:

\begin{lstlisting}[style=py]
est = LCMCovariance(dim=X.shape[1], dtype=torch.float64)
est.fit_moments(X)                      # natural index order
est.fit_moments(X, order=my_order)      # or supply an ordering
nll = est.nll(X_heldout)                # O(dim) Kalman recursion, no dense matrix
\end{lstlisting}

\paragraph{Provenance.} The ingredients are classical: this is a Yule--Walker-style moment
estimator for an OU/AR(1) process. What is not textbook is the \emph{non-stationary} case ---
a separate amplitude per coordinate, a separate decay per gap, and a per-coordinate nugget,
where standard treatments assume one variance and one decay. Two targeted literature searches
did not turn it up, but two searches are not a review, so \textbf{novelty should be treated as
unverified}. What is verified is that it is absent from the paper: that paper fits by
minimising a squared Frobenius distance, derived in its own appendix, and its released code
contains no moment estimator, no band-based fitting and no closed form of any kind.

\section{Relationship to the published method}
\label{sec:relation}

The paper minimises $\lVert \Sigma - \hat\Sigma \rVert_F^2$ over $(d, w, a)$ jointly with Adam
--- learning rate $0.01$, 200 epochs in its channel configuration, 50 in its inversion
configuration. The two procedures do \emph{not} solve the same problem.

\begin{center}
\begin{tabular}{@{}lll@{}}
\toprule
 & Paper's gradient fit & Moment matching \\
\midrule
Ordering of the variables & \textbf{learned} (sorting by $a$) & \textbf{given} (an input) \\
Coordinate spacing        & learned                            & closed form \\
Amplitudes and nugget     & learned                            & closed form \\
Objective                 & squared Frobenius distance         & moment identity, none \\
Cost                      & iterative, tuned                   & one pass, $O(N)$ \\
\bottomrule
\end{tabular}
\end{center}

So moment matching addresses a \textbf{strictly smaller problem}: the parameters, conditional
on an ordering. The paper's method is more ambitious and attempts both at once.

\subsection{Why comparing them on flattened feature maps is still fair}

The paper initialises $a = \operatorname{linspace}(-1, 1, N)$, which assigns coordinate $i$ in
proportion to $i$ --- so it already encodes the natural index order. On a flattened feature map
that order is meaningful, since neighbouring indices are adjacent pixels. Both procedures
therefore start from the same ordering information, and what differs between them is the
\emph{spacing}.

Table~\ref{tab:exp25} gives the held-out Gaussian cost of centred CIFAR-100 features, in nats
per dimension, lower better, 3 seeds, at dimensions up to $196\,608$.

\begin{table}[htbp]
\centering
\caption{Fitting rules compared at the paper's own scale. Lower is better.}
\label{tab:exp25}
\begin{tabular}{@{}lcccc@{}}
\toprule
Fitting rule & par/dim & $16\,384$ & $65\,536$ & $196\,608$ \\
\midrule
Diagonal --- the paper's own baseline          & 1 & $0.8372$ & $0.7566$ & $1.0030$ \\
Factor analysis, 3 factors (larger budget)     & 4 & $0.6669$ & $0.6155$ & $0.8425$ \\
Chain, exact maximum likelihood                & 3 & $0.6811$ & ---      & ---      \\
Chain, the paper's Frobenius rule              & 3 & $0.6000$ & $0.6165$ & $0.8472$ \\
Chain, Frobenius, coordinates \emph{frozen}    & 2 & $0.6349$ & $0.5969$ & $0.8396$ \\
\textbf{Chain, moment matching}                & 3 & $\mathbf{0.1325}$ & $\mathbf{-0.0537}$ & $\mathbf{-0.2873}$ \\
\bottomrule
\end{tabular}
\end{table}

Paired on the same seeds, the advantage over their rule is $+0.4674 \pm 0.0508$ at $16\,384$
(9 standard errors), $+0.6702 \pm 0.0095$ at $65\,536$ (71), and $+1.1345 \pm 0.0229$ at
$196\,608$ (50). Fitting time is $98$\,s against $1723$\,s at the largest width.

\subsection{What that comparison shows, and what it does not}

\paragraph{Their fitting is the weak link, not their loss.} Fitting the same family by exact
maximum likelihood --- optimising precisely the quantity being scored, at $180\times$ the cost
--- reaches $0.6811$, slightly \emph{worse} than their Frobenius rule. Both gradient routes
stall in the same region, so the difficulty is the optimisation landscape of this
parameterisation rather than the choice of objective.

\paragraph{The mechanism is the coordinate span.} Expressing decay over a few neighbours at $N$
coordinates needs a span of order $N$. The published initialisation gives a span of $2$
whatever the dimension, which at $N = 200\,704$ puts adjacent coordinates at correlation
$0.99999$ and the most distant pair at $e^{-2} = 0.135$: almost structureless, and roughly
$400\,000\times$ away from the solution. Moment matching does not travel there, because
$\operatorname{cumsum}(-\log\varphi)$ \emph{computes} the span. Table~\ref{tab:span} shows what
each fit produces on a ResNet-34's four layers.

\begin{table}[htbp]
\centering
\caption{What each fitted covariance actually imposes, per layer.}
\label{tab:span}
\begin{tabular}{@{}lccc@{}cccc@{}}
\toprule
& \multicolumn{3}{c}{Paper's Frobenius fit} & & \multicolumn{3}{c}{Moment matching} \\
\cmidrule(r){2-4}\cmidrule(l){6-8}
Layer & span & nbr.\ corr. & fit/real var. & & span & nbr.\ corr. & fit/real var. \\
\midrule
layer1 & $\mathbf{0.0}$ & $0.0002$ & $3.06\times$  & & $817\,775$ & $0.569$ & $1.000$ \\
layer2 & $\mathbf{0.0}$ & $0.0000$ & $6.81\times$  & & $340\,766$ & $0.443$ & $1.000$ \\
layer3 & $\mathbf{0.0}$ & $0.0000$ & $16.03\times$ & & $141\,889$ & $0.477$ & $1.000$ \\
layer4 & $1.2$          & $0.0002$ & $0.63\times$  & & $16\,476$  & $0.849$ & $1.000$ \\
\bottomrule
\end{tabular}
\end{table}

With every coordinate coincident $K(a)$ is all-ones, so $\Sigma$ can only be
diagonal-plus-rank-one --- and the rank-one weights vanish too, which is what the $0.0002$
measures. What is left is a diagonal with variances $3$ to $16\times$ the real feature
variances. The moment fit matches those variances exactly, by construction, and imposes real
neighbour correlation of $0.44$ to $0.85$.

\paragraph{A second, more direct cause: their optimiser decays the coordinate.} Their
inversion script sets the parameter groups as

\begin{lstlisting}[style=py]
# reference/REMIX/efficient_inversion/resnet34/synthetic_samples.py:117-120
optim.Adam([{'params': lcms[l].u_raw, 'lr': 1e-2},
            {'params': lcms[l].w,     'lr': 1e-2},
            {'params': lcms[l].a,     'lr': 0.12, 'weight_decay': 1e-5}],
           betas=(0.9, 0.999))
\end{lstlisting}

\noindent
Weight decay on a \emph{coordinate} is a pull toward the origin --- that is, toward a collapsed
span and an all-ones kernel. The two causes compose: at span $2$ and $N = 200\,704$ the kernel
is within $10^{-5}$ of all-ones, so the data gradient on $a$ is negligible, and the decay term
is then the only consistent signal acting on it. Over 50 epochs at learning rate $0.12$ that is
ample to drive the span to zero. Tellingly, their \emph{likelihood} script --- the one producing
their Figure~4 --- uses the same learning rate for $a$ with \textbf{no} weight decay
(\code{calculate\_likelihood.py:165}). The decay appears only in the image-generation path,
which is exactly where we measured the collapse.

\paragraph{It shows nothing about their ordering claim.} Both arms used the natural index. Note
also the frozen-coordinate row of Table~\ref{tab:exp25}: fixing the coordinates at equal
spacing uses one parameter less and beats \emph{learning} them at both large widths. On that
evidence the learned ordering is not earning its keep --- but the evidence is itself from
gradient fits starting at span $2$, so it bounds the published procedure rather than the idea
of learning an order.

\section{The ordering requirement, measured}
\label{sec:ordering}

Because the bands are the chain's neighbours only when the coordinates are already sorted, a
wrong order is not a small perturbation. Table~\ref{tab:order} reports held-out negative
log-likelihood on data drawn from a known chain, 3 seeds. ``Order recovery'' is the absolute
correlation between a spectrally seriated order and the true one: $1.0$ perfect, $0$ chance.
The diagonal column is the no-structure floor.

\begin{table}[htbp]
\centering
\caption{Sensitivity of the closed-form fit to the ordering it is given. Lower is better.}
\label{tab:order}
\begin{tabular}{@{}rrrrrrr@{}}
\toprule
$N$ & adj.\ corr. & order recovery & true order & seriated & random & diagonal \\
\midrule
$256$  & $0.575$ & $\mathbf{0.999}$ & $361.7$  & $370.4$  & $434.3$  & $437.0$  \\
$256$  & $0.271$ & $0.270$          & $426.8$  & $436.6$  & $436.9$  & $437.0$  \\
$1024$ & $0.575$ & $0.690$          & $1446.3$ & $1674.5$ & $1748.4$ & $1750.2$ \\
$2048$ & $0.575$ & $0.531$          & $2892.3$ & $3441.7$ & $3497.6$ & $3498.8$ \\
$2048$ & $0.271$ & $0.019$          & $3418.5$ & $3499.7$ & $3499.8$ & $3499.8$ \\
\bottomrule
\end{tabular}
\end{table}

Two readings. \textbf{A random order yields a diagonal}, landing within $1.3$ nats of the floor
at every setting: the implied amplitudes go to zero and the fit collapses, with no graceful
degradation. And \textbf{spectral seriation substitutes for the true order only at small
dimension and strong coupling} --- essentially exact at $N = 256$ with adjacent correlation
$0.575$, falling to $0.531$ at $N = 2048$ and to $0.019$ once adjacent correlation drops to
$0.27$, with the fit degrading alongside it.

So where the data supplies a meaningful order --- a flattened feature map, a spectrum, a time
series --- moment matching is immediately usable. Where the order must be discovered, there is
no demonstrated way to supply it at the dimensions this family targets.

\section{Where it does not work}
\label{sec:limits}

\paragraph{It emits only positive couplings.} Taking $w = +\sqrt{w^2}$ discards the sign, so
every off-diagonal entry it produces carries the sign of $w_p w_q$. This is harmless where
neighbours are positively correlated, and fatal over \emph{channels}, where $47$--$49\%$ of
off-diagonal entries are negative: the estimator collapses towards the identity, scoring
$0.986$--$0.995$ relative off-diagonal Frobenius error where $1.0$ means predicting no
correlation at all.

The sign is in principle recoverable from the same bands at no parameter cost, since the decay
is positive and therefore $\operatorname{sign}(\Sigma_{p,p+1}) = \operatorname{sign}(w_p
w_{p+1})$ propagates along the chain. It does not survive testing: a small-scale check improved
a seriated fit from $0.98$ to $0.88$, but across 3 seeds $\times$ 4 layers it reached $0.9897$,
no better than unsigned.

\paragraph{It uses three bands only}, so it cannot exploit longer-range structure and has no
recourse when the target is poorly described by a chain at all.

\paragraph{A better fit is not automatically better downstream.} Dropping this covariance into
the paper's own inversion loop made the generated images \emph{worse}: correlation-profile
error $0.2184$ against $0.0623$ for their own fit and $0.0250$ for their no-covariance
ablation, and Fr\'echet distance $368.7$ against $310.7$. The reason is structural rather than a
bug. Minimising a Gaussian negative log-likelihood \emph{with respect to the data} is
mode-seeking, so a correctly correlated precision penalises differences between neighbouring
coordinates and smooths the feature map. Better density estimation, stronger smoothing prior.

\section{How this could improve the published method}
\label{sec:improve}

All six proposals below are \textbf{untested}. They are ordered cheapest first, and each says
what would falsify it. They apply to the flattened-feature configuration; see the caution in
Section~\ref{sec:caution} about the channel configuration.

\subsection{Remove the weight decay on the coordinate --- one character}

Their inversion script decays $a$ toward zero (\code{synthetic\_samples.py:119},
\code{weight\_decay: 1e-5}), which is a pull toward exactly the collapsed, all-ones
configuration measured in Table~\ref{tab:span}. Their own likelihood script omits it. Deleting
it from the inversion path costs nothing and makes the two consistent.

\emph{Falsified if} the span still collapses without it --- which would mean the vanishing
data gradient at span $2$ suffices on its own, and the next proposal is the real fix. Either
outcome identifies the cause, which is why this is the first thing to run.

\subsection{Scale the coordinate initialisation with dimension --- one line}

The span of $2$ is dimension-independent, while the required span grows with $N$. Replace
$\operatorname{linspace}(-1,1,N)$ with a span set from the data: the mean absolute band-1
correlation $\rho_1$ implies a gap of $-\log \rho_1$, so initialise
$a = \operatorname{arange}(N)\cdot(-\log\rho_1)$. This needs only $s_1$, costs one pass, and
leaves their objective, optimiser and code path untouched.

\emph{Falsified if} the fit still collapses the span, or ends no better than from
$\operatorname{linspace}(-1,1,N)$. A caution from our own data: over channels we tried a span of
$N$ and the fit became substantially worse, because channel covariance is broad rather than
local. The right span is whatever the data implies, not simply a large one.

\subsection{Warm-start their gradient fit from the closed form}

Run moment matching, then run their Frobenius descent from the result. This keeps their
objective and their ability to refine beyond what three bands identify --- longer-range
structure, the two unidentified endpoints --- while starting in the region the optimiser cannot
reach unaided.

\emph{Falsified if} their descent moves away from the moment solution and ends worse than it.
That outcome would itself be informative: it would mean the Frobenius optimum lies genuinely
elsewhere, and the gap in Table~\ref{tab:exp25} is a disagreement between objectives rather
than an optimisation failure.

\subsection{Alternate closed-form blocks with gradient steps on the ordering only}
\label{sec:alternate}

This is the proposal the structure of the problem most directly supports. The parameters split
cleanly by what each method is good at:
\begin{itemize}
  \item Given the ordering and the decays, \textbf{amplitudes and nugget have closed-form
        solutions}.
  \item Given the ordering and the amplitudes, \textbf{the decays have a closed-form solution}
        from band 1.
  \item The \textbf{ordering} has no closed form at all --- and it is precisely what their
        gradient fit is for.
\end{itemize}

\begin{lstlisting}[style=py,caption={Alternating scheme: closed form for scale, gradients for order.},label={lst:alt},float=htbp]
order = natural_index            # or a seriation of |corr|, or their linspace ordering
for sweep in range(max_sweeps):
    # 1. closed form: reset scale parameters to their optimum for THIS ordering.  O(N)
    a, w, d = fit_moments(X, order=order)
    model.a.data, model.w.data, model.d.data = a, w, d

    # 2. gradient steps on the coordinates ONLY, with the scales frozen.
    #    Their objective, their optimiser -- just a much smaller parameter set.
    opt = torch.optim.Adam([model.a], lr=lr)
    for _ in range(inner_steps):
        opt.zero_grad()
        frobenius_loss(model.covariance(), S_hat).backward()   # their Eq. (5)
        opt.step()

    # 3. the ordering is the discrete state; stop when it stops moving.
    new_order = torch.argsort(model.a)
    if torch.equal(new_order, order):
        break
    order = new_order
\end{lstlisting}

Step 1 keeps the scale parameters at their moment-matched optimum for the current ordering, so
the gradient in step 2 never has to fix the span --- it only has to decide the permutation,
which is the one thing moments cannot give. Step 3 makes the alternation a discrete descent on
the ordering, with a natural stopping rule.

Two properties make this attractive. The expensive part of their fit, iterating on all $3N$
parameters, shrinks to iterating on $N$. And the silent failure mode disappears: the span
cannot collapse, because step 1 resets it from the data every sweep.

\emph{Falsified if} the alternation does not improve on moment matching with a fixed natural
order, or if it cycles without converging. The honest prior is mixed. Section~\ref{sec:ordering}
shows the ordering is worth a great deal when you have it, and gradient fits in this project
have never recovered an ordering from an uninformed start --- but all of those fits also began
at the wrong span, which this scheme removes. That confound is exactly what makes the
experiment worth running.

\subsection{Use the closed-form span as a free diagnostic}

The collapse in Table~\ref{tab:span} is invisible in the loss curve, which falls smoothly and
looks converged while the model degenerates to a diagonal. Computing the moment-matched span
costs one pass and gives a reference: if the fitted span is orders of magnitude smaller, the
fit has collapsed. This catches the failure without changing their method at all, and is the
one proposal here with essentially no downside.

\subsection{For channels, address the sign pattern rather than the fitting}

Over channels the limitation is representational, not numerical. The achievable sign pattern is
$\operatorname{sign}(w_i w_j)$, which is rank-one --- a two-camp split --- while almost half the
real off-diagonal entries are negative in an overlapping pattern. A principled attempt would
estimate the camp assignment separately, for instance from the sign of the leading eigenvector
of the empirical correlation, then moment-match within camps.

This is the most speculative proposal here, and the evidence is against it: the simpler sign
restoration already failed at scale, and over channels a rank-two factor model at one extra
parameter per channel beats the chain's best achievable fit by a wide margin. The honest
conclusion over channels is that fitting is not where the problem lies.

\subsection{The caution that governs all of these}
\label{sec:caution}

\textbf{Improving the fit may not improve the paper's results.} Section~\ref{sec:limits} records
that a better-fitting covariance made their inversion output measurably worse, because the
fitted covariance enters their loss as a penalty on the generated features and a correctly
correlated penalty over-smooths. Anyone pursuing the first three proposals should therefore
measure the downstream metric, not only fit quality, and should expect the covariance weight to
need re-tuning: their loop multiplies this term by a single fixed weight, and a better-fitted
covariance hands it a term of quite different scale and curvature.

There is a reading of these results under which the paper's gradient fit is \emph{accidentally}
doing something useful: collapsing to a loose, over-wide diagonal that acts as a mild
regulariser. If that is what drives the reported gains, then fixing the fit removes the
mechanism. Establishing which is true needs the one experiment nobody has run --- turning the
covariance term off and reading the accuracy.

\section{Protocols}
\label{sec:protocols}

Every number quoted above comes from one of the three experiments below, each given in enough
detail to re-run without reading the code.

\subsection{The fitting comparison at scale (Table~\ref{tab:exp25})}

\paragraph{Data.} CIFAR-100, 3000 images from the training split with a per-seed shuffle.
Images are used only to produce features; labels only to train the network.

\paragraph{Network.} A deliberately stride-1 convolutional stack, so the flattened maps reach
the paper's dimensions: three blocks of
\code{Conv2d(3$\times$3, padding 1)} $\to$ \code{BatchNorm2d} $\to$ \code{ReLU} with 16, 64 and
192 channels, then global average pooling into a linear head. \textbf{The network is trained},
not random --- 4 epochs of Adam, learning rate $2\mathrm{e}{-3}$, batch 64, cross-entropy.

\paragraph{Features.} The flattened output of each block, so $D = C \times 32 \times 32$ gives
$16\,384$, $65\,536$ and $196\,608$. Split $80/20$ into fitting and held-out sets, centred on
the \emph{fitting} mean.

\paragraph{Cleaning, applied identically to every arm.} ReLU maps contain near-dead units whose
held-out cost is decided purely by whatever floor each family clamps at --- at $D = 65\,536$,
seven units have variance below $10^{-6}\times$ the median and one is exactly zero. Those
dimensions are dropped (variance below $10^{-6} \times$ median), and every family is then given
the \emph{same} explicit variance floor of $10^{-4} \times$ median. Before this was fixed the
diagonal scored $35.8$ nats/dim better purely on floor choice.

\paragraph{Arms.}

\begin{center}
\begin{tabular}{@{}lll@{}}
\toprule
Arm & Params & How it is fitted \\
\midrule
Diagonal & $D$ & per-dimension variance of the fitting split, clamped \\
Factor analysis, rank 2 / 3 & $(r{+}1)D$ & EM, 40 iterations, randomised-SVD init \\
Chain, Frobenius (paper's rule) & $3D$ & Adam, 300 steps, lr $5\mathrm{e}{-2}$, minibatch 256 \\
Chain, Frobenius, $a$ frozen & $2D$ & identical, $a$ fixed at equal spacing \\
Chain, exact maximum likelihood & $3D$ & Adam on $(a,w,u)$, 100 steps, lr $5\mathrm{e}{-2}$ \\
\textbf{Chain, moment matching} & $3D$ & closed form, one pass, no optimiser \\
\bottomrule
\end{tabular}
\end{center}

\noindent
The exact-likelihood arm is skipped above $D = 65\,536$: it needs ${\sim}1150$\,s even at
$16\,384$.

\paragraph{Metric.} Held-out Gaussian negative log-likelihood, summed over dimensions, averaged
over held-out images, divided by $D$ --- hence \emph{nats per dimension}, lower better.
Negative values are legitimate differential entropies, and are why the result was
cross-checked against a dense Cholesky.

\paragraph{Seeds and cost.} 3 seeds, re-seeding the network initialisation, the image sample and
the fitting. About 5\,h on one CPU node of the GMUM cluster, 10 cores, 48\,GB, float64
throughout.

\subsection{What each fit imposes, per layer (Table~\ref{tab:span})}

\paragraph{Data.} 256 real ImageNet images of one class (\code{n03457902}, greenhouse) at
$224 \times 224$.

\paragraph{Network.} Torchvision ResNet-34 with ImageNet weights, frozen, in eval mode.
Features are taken at the outputs of \code{layer1}--\code{layer4} and flattened, giving $D$ up
to $200\,704$.

\paragraph{Fitting.} Their objective with \textbf{their own optimiser settings}, imported from
their clone rather than reimplemented: Adam with per-group learning rates \code{u\_raw}
$10^{-2}$, \code{w} $10^{-2}$, \code{a} $0.12$ with \code{weight\_decay} $10^{-5}$,
\code{betas} $(0.9, 0.999)$, 50 epochs over the 256 images at batch 32. The moment arm instead
fits in closed form from the same features.

\paragraph{What is reported.} The span of the fitted coordinates, the median correlation the
fit implies between adjacent coordinates, and the fitted variance divided by the real feature
variance --- none of which is visible in the loss curve.

\paragraph{Cost.} Minutes: the inversion loop is truncated, and the fit precedes it and does
not depend on it. Cluster job 542802, one 24\,GB card.

\subsection{Ordering dependence (Table~\ref{tab:order})}

\paragraph{Data.} Synthetic, drawn from a known chain in the family --- the only setting where
a true ordering exists to compare against. Coordinates evenly spaced at gap $g$, amplitudes
drawn uniformly from $[0.8, 1.3]$, 4000 samples for fitting and 4000 held out.

\paragraph{Sweep.} Dimension $256/1024/2048$ crossed with gap $0.05/0.223/0.8$, giving adjacent
correlations of roughly $0.58/0.48/0.27$. 3 seeds each.

\paragraph{Arms.} The same closed-form fit given four orderings: the true one; a spectral
seriation computed from the data alone (Fiedler vector of the Laplacian of the absolute
sample-correlation graph); a uniform random permutation; and an independent-coordinates
Gaussian as the no-structure floor.

\paragraph{Metric.} Held-out Gaussian negative log-likelihood summed over dimensions and
averaged over images --- \emph{not} divided by $D$ here, so columns are comparable only within
a row. ``Order recovery'' is the absolute Pearson correlation between the seriated order and
the true index, measuring monotonicity up to reflection, which is what the chain needs.

\paragraph{Cost.} About 3 minutes on CPU.

\section{Reproducing what is quoted here}

\begin{lstlisting}[style=py,language=bash,basicstyle=\ttfamily\footnotesize]
# fitting comparison at the paper's scale (cluster, ~5 h on one CPU node)
sbatch cluster/exp25.sbatch

# correctness against a dense Cholesky (laptop, seconds)
python -u experiments/exp25_verify_moments.py

# ordering dependence and seriation's reach (laptop, ~3 min on CPU)
python -u experiments/exp30_ordering_dependence.py
\end{lstlisting}

\begin{center}
\begin{tabular}{@{}p{0.33\textwidth}p{0.59\textwidth}@{}}
\toprule
What & Where \\
\midrule
The estimator & \code{Laplace/laplace/lcm.py} --- \code{LCMCovariance.fit\_moments} \\
Fitting comparison at scale & \code{experiments/exp25\_paper\_scale\_baselines.py} \\
Correctness check & \code{experiments/exp25\_verify\_moments.py} \\
Ordering dependence & \code{experiments/exp30\_ordering\_dependence.py} \\
Span collapse in the paper's loop & \code{experiments/exp28\_model\_inversion.py} \\
The signed variant that failed & \code{experiments/exp29\_channel\_correlation\_fit.py} \\
The paper's code, cloned & \code{reference/REMIX/} \\
\bottomrule
\end{tabular}
\end{center}

\noindent
Every result file has a \code{.meta.json} beside it recording the configuration, seeds, both
git commits, host and, for cluster runs, the job id.

\end{document}
